\subsection{Mahindra University: history and fundings}

Mahindra École Centrale College of Engineering was founded in 2014 in Hyderabad,
Telangana, India, by the École Centrale group. It is sponsored by Mahindra
Educational Institutions, a subsidiary of Tech Mahindra, from the Mahindra
Group in India. In 2020 it became part of the Mahindra University (MU), not only
having a school of Engineering but also of Management, Law, Education, Digital
Media and Communication, Design Innovation and Hospitality Management. 

The Mahindra group owns multiple businesses and is very famous in India for its
role in automotive, the Mahindra brand makes cars, motorcycles, buses, trucks,
tractors. The group also plays a role in finance, IT, logistics, real estate,
and energies. To give an idea of the size of the Mahindra group, it counts 324k
employees.

The Mahindra University Engineering school offers 15 B.Tech options (Bachelor
degree in technology) and 10 M.Tech options (Master degree in technology) in the
following fields: Physics, Civil Engineering, Electrical and Computer
Engineering, Chemistry, Mathematics, Mechanical and Aerospace Engineering,
Computer Science and Engineering, Entrepreneurship, Humanities and Social
Sciences, Life Sciences, Natural Sciences. They also have many PhD in
Engineering which they deliver in 4 years.

It is located in Jeedimetla in the northern outskirt of Hyderabad, Telangana,
next to the Tech Machindra company, but far from other universities research
work are in collaboration with, like University of Hyderabad (UoH), IIT
Hyderabad or IIIT Hyderabad, see \cref{fig:mu-map}.

\insertfigure
    {assets/images/MU-map}
    {width=0.7\textwidth}
    {MU on the map of Hyderabad (map from OpenStreetMap)}
    {fig:mu-map}

\subsection{Research teams and projects}

During my internship I assisted Dr Nidhi Goyal, my supervisor, in many tasks. Dr
Nidhi Goyal is Assistant Professor in the Computer Science and Engineering
departement of MU, she specializes her research in Natural Language Processing,
Knowledge Graphs, Machine Learning, and their applications to Neuroscience.
 
The team I was in was composed of Dr Nidhi Goyal, Dr Sai Amrit Patnaik (AI
Researcher), Gaspard Lefort (Intern), Shankar Sai Ankit Reddy (B.Tech Student)
and me. It aimed to publish papers on generative AI and EEG, one paper on EEG to
Image generative models (mine) and the other on EEG to Video generative models
(Gaspard's). When Shankar Sai Amrit Patnaik was working from Delhi, the rest of
the team was in MU, so most of the meetings were from planned videocalls.  As Dr
Nidhi Goyal is Assistant Professor, she couldn't be working full time on the
research projects, she was also helping with clubs in the university, but when
the summer vacations for students started, she could be working full time on
research projects.

Then, Gaspard and I helped to set up experiments about attention of car drivers,
subjects drove on a driving simulator and were subject to disturbances (phone
notification, dashboard warning light, storm), an eye tracker and a EEG headset
monitored their reactions. The project was established through a collaboration
between Mahindra University (MU) and the University of Hyderabad (UoH), the team
was composed of Dr Sunder Bukya (MU), Dr Nidhi Goyal (MU), Dr Shiva Ram Reddy
(UoH) and Srujan Kumar (B. Tech student from UoH).

\subsection{Differences with French universites}

In MU there were many differences with most universities in France, part of it
is because India is culturally very different from France but also because
Mahindra University is a private school.

Students in the University, but also PhDs and part of the researchers and
teachers stay in the campus' hostel, bachelor and master students require
permission to leave the campus. Students had many clubs in the university, and
many of these clubs were study-related, students faced a high competition to be
part of some clubs, they had to apply and go through an interview.

There were also many cultural events about litterature, spirituality or the
"francophony week", where they showed French movies or French food.  Bachelor
students had a mantaory French language course as part of the Ecole Centrale
group connection. 

\subsection{Using a supercomputer (NVidia DGX-A100)}

MU has a NVidia DGX-A100 Cluster contains 8 DGX-A100 GPUs with 64 CPU cores. It
is mainly used for machine learning tasks are they require a lot of parallel
computational capabilities.

Using this supercomputer showed many difficulties.
\begin{itemize}
\item The memory limit
Each account on the supercomputer had a limit of 100Go which was not enough to
carry all the needed data. Datasets could weigh around 30Go, many checkpoints
were around 5 to 10Go large, miniconda used around 25Go for all the
environments we could not to mix for reproducibility purposes.

\item CUDA learning
The cuda related commands had to be learned, to launch computations, to grant
enough privileges, to allocate enough GPU VRAM.

\item Constrained environment
With the DGX-A100 cluster it was not possible to use Docker containers, and the
installed CUDA version was too old to run some codes, in order to have more
recent CUDA versions it was needed to use a NVidia CUDA container with the
enroot command.

\item Availability of the cluster
Even though the cluster had two nodes, it could not handle too many allocations
in the same time, making it impossible to use the cluster half of the time
because too many people were using it. 
\end{itemize}

The supercomputer system admin, Aman Iftekhar, was always very easy to reach, he
helped a lot through the internship and teached helpful commands to use the
cluster.